\documentclass[aspectratio=169]{beamer}
\usetheme{Madrid}
\usepackage{amsmath}
\usepackage{booktabs}

\title{Progress Update: Irregular Time-Series Autoencoder}
\author{Haochen Zhang}
\date{\today}

\begin{document}

\begin{frame}
  \titlepage
\end{frame}

\begin{frame}{What I Am Building}
  \begin{itemize}
    \item Goal: learn a modality-specific autoencoder for \textbf{irregular AI-READI daily windows}.
    \item Current focus file: \texttt{train\_irregular\_ts\_ae.py}.
    \item Model file: \texttt{models/modality\_imputation/ts\_ae.py}.
    \item Target use: robust reconstruction/imputation for wearable modalities (e.g., heart rate, respiratory rate, physical activity, calorie), optionally aligned to glucose timeline.
  \end{itemize}
\end{frame}

\begin{frame}{Data Construction Pipeline}
  \begin{itemize}
    \item Dataset class: \texttt{AIREADIModalityImputationDataset} with daily windows.
    \item Typical setup:
    \begin{itemize}
      \item anchor modality: glucose
      \item sequence length: 288 (5-min bins over 24h)
      \item aligned modality option: \texttt{<modality>\_aligned\_to\_glucose}
    \end{itemize}
    \item I cache and filter samples in \texttt{IrregularTSDataContainer}:
    \begin{itemize}
      \item keep only samples with valid \((T, C)\) value/mask pairs
      \item remove windows above \texttt{max\_missing\_ratio}
    \end{itemize}
    \item Output to model:
    \begin{itemize}
      \item \(\mathbf{x}\in\mathbb{R}^{T\times C}\): values
      \item \(\mathbf{m}_{obs}\in\{0,1\}^{T\times C}\): observed mask
    \end{itemize}
  \end{itemize}
\end{frame}

\begin{frame}{Model Architecture (\texttt{IrregularTimeSeriesAE})}
  \begin{itemize}
    \item Input projection: linear \(C \rightarrow d_{model}\)
    \item Positional encoding: sinusoidal
    \item Core encoder: Transformer encoder (multi-layer, multi-head self-attention)
    \item Output projection: linear \(d_{model} \rightarrow C\)
    \item Key idea: one learnable \texttt{special\_token} replaces:
    \begin{itemize}
      \item dataset-missing/padded timesteps
      \item additionally masked timesteps (for denoising-style training)
    \end{itemize}
    \item No attention mask is applied, so all timesteps can attend globally.
  \end{itemize}
\end{frame}

\begin{frame}{Training Objective and Mask Strategy}
  \textbf{Extra masking during training}
  \begin{itemize}
    \item On observed entries only, randomly drop a fraction (\texttt{input\_random\_drop\_prob}).
    \item This creates \(\mathbf{m}_{in}\in\{0,1\}^{T\times C}\), where 1 means ``hidden from input''.
  \end{itemize}

  \textbf{Loss}
  \[
    \mathcal{L}
    =
    \frac{\sum (\hat{\mathbf{x}}-\mathbf{x})^2 \odot \mathbf{m}_{in}}
         {\max\left(1,\sum \mathbf{m}_{in}\right)}
  \]

  \begin{itemize}
    \item So the model is optimized specifically to recover the additionally hidden observed values.
    \item Validation uses the same masking policy to measure reconstruction under controlled sparsification.
  \end{itemize}
\end{frame}

\begin{frame}{Current Training Setup in Practice}
  \begin{itemize}
    \item Script: \texttt{scripts/train\_irtsae.sh}
    \item Typical hyperparameters:
  \end{itemize}

  \begin{table}
    \centering
    \begin{tabular}{ll}
      \toprule
      Setting & Value \\
      \midrule
      \(d_{model}\) / heads / layers & 128 / 4 / 3 \\
      batch size & 256 \\
      optimizer & AdamW \\
      learning rate & \(3\times10^{-4}\) \\
      epochs & 1000 \\
      random drop prob & 0.3 \\
      max missing ratio & 0.5 \\
      \bottomrule
    \end{tabular}
  \end{table}

  \begin{itemize}
    \item Best checkpoint is selected by validation masked-MSE and saved as \texttt{best.pt}.
    \item Reconstruction plots are generated periodically for qualitative checks.
  \end{itemize}
\end{frame}

\begin{frame}{What I Will Report Next}
  \begin{itemize}
    \item Quantitative:
    \begin{itemize}
      \item validation masked-MSE by modality
      \item sensitivity to \texttt{input\_random\_drop\_prob} and \texttt{max\_missing\_ratio}
    \end{itemize}
    \item Qualitative:
    \begin{itemize}
      \item side-by-side reconstructions with/without extra masking
      \item failure cases at highly sparse windows
    \end{itemize}
    \item Planned extension:
    \begin{itemize}
      \item compare against simple interpolation baselines
      \item test whether latent representations help downstream drift modeling
    \end{itemize}
  \end{itemize}
\end{frame}

\end{document}
